\section{Componentes necesarios del Transformer Block}

Escribir solo la attention equation deja fuera la mayor parte del modelo real. El orden de los tokens, la optimización en profundidad, la transformación no lineal y los límites de la información quedan a cargo, respectivamente, de la positional representation, de residual/normalization, de la feed-forward network y de la mask. No son adornos: son la base que permite que attention sea entrenable, apilable y aplicable a distintas tareas.

\subsection{Position, residual, normalization y masking}

La self-attention pura tiene permutation equivariance respecto del orden de la entrada: si se reordenan los tokens a la vez, la salida se reordena de la misma manera, pero el modelo no conoce el significado absoluto de «el primero» y «el último». La \term{positional encoding} (codificación posicional) inyecta el orden en la token representation. El Transformer original usa seno y coseno:
\[
PE(pos,2i)=\sin\!\left(pos/10000^{2i/d}\right),\qquad
PE(pos,2i+1)=\cos\!\left(pos/10000^{2i/d}\right).
\]
$pos$ es la posición, $i$ es el índice del canal y $d$ es la dimensión del modelo. Los modelos modernos también usan con frecuencia position learned, relative o rotary, pero el objetivo común es expresar el orden relativo o absoluto.

\lecturefigure{slide-13-ingredients.jpg}{Position, normalization, residual connection y masking forman juntos un Transformer block utilizable.}{Video oficial, 12:30; Vaswani et al. 2017.}

\begin{knowledgebox}{Lectura de la figura: cada ingredient corrige un modo de falla distinto}
Position corrige la ausencia de orden; la residual connection permite que la capa aprenda incrementos y ofrece un atajo para el gradiente; la layer normalization controla la escala de cada token representation; la feed-forward network agrega, en cada posición de forma independiente, una transformación no lineal entre canales; la mask establece qué información es legítimamente visible. Si se elimina cualquiera de estos componentes, el modelo quizá siga funcionando, pero ya no tendrá la estabilidad de entrenamiento ni la semántica de tarea de la arquitectura original.
\end{knowledgebox}

La \term{residual connection} (conexión residual) suele escribirse como $y=x+F(x)$, de modo que el block aprende una corrección respecto de la entrada. La \term{layer normalization} (normalización por capa) estandariza la dimensión de features de cada token y facilita la optimización de redes profundas. La position-wise feed-forward network es
\[
\operatorname{FFN}(x)=W_2\sigma(W_1x+b_1)+b_2,
\]
que comparte parámetros entre todas las posiciones pero calcula cada una de forma independiente: attention se encarga del token mixing y la FFN del channel mixing.

\begin{warningbox}{La mask es un límite semántico, no solo un parámetro de rendimiento}
Si un causal language model omite la causal mask, durante el entrenamiento verá la información futura a la derecha del token objetivo; la loss puede verse anormalmente buena, pero la generación resultará completamente inútil. Si la padding mask es incorrecta, los símbolos de relleno se tratarán como contenido real. Cualquier implementación de attention debe probar primero la matriz de visibilidad y después el throughput.
\end{warningbox}

\subsection{Resumen de la sección}

El Transformer block combina token mixing y channel mixing, y mediante position, mask, residual y normalization conserva el orden, la causalidad y la capacidad de entrenamiento. Recordar esta división de funciones es más útil que memorizar un diagrama de arquitectura.
